\subsection{Memory-Kernel Origin of the SBC/F3 Infrared Sector}

The observationally viable SBC/F3 branch suggests that the dominant correction is not a homogeneous deformation of the background expansion, but an effective infrared response of the perturbative sector.

We model this response through a causal memory kernel,
\begin{equation}
\Delta_X(z)=\int_z^{\infty} K(z,z')\,\mathcal{R}_X(z')\,dz',
\end{equation}
where \(X\) denotes a cosmological observable, \(K(z,z')\) is a causal memory kernel, and \(\mathcal{R}_X\) is the perturbative response source.

A minimal exponentially screened kernel is
\begin{equation}
K(z,z') = K_0 \exp[-k_0 |z'-z|],
\end{equation}
with \(k_0>0\). For slowly varying infrared sources, the leading effective response reduces to
\begin{equation}
F(z)\simeq -\epsilon \exp[-k_0(1+z)].
\end{equation}

This corresponds to the SBC/F3 phenomenological form
\begin{equation}
F(z)= -\epsilon
\exp\left[-\left(k_0(1+z)\right)^n\right],
\end{equation}
with the leading memory-decay regime
\begin{equation}
n\simeq 1.
\end{equation}

The repeated observational preference for \(n\simeq 1\) can therefore be interpreted as evidence that the viable SBC/F3 sector behaves as a first-order exponentially screened infrared memory response.

The parameter \(k_0\) controls the characteristic infrared screening scale, while \(\epsilon\) fixes the amplitude of the effective perturbative response. The observationally preferred values,
\begin{equation}
\epsilon\simeq 0.06,\qquad
k_0\simeq 0.015-0.022,\qquad
n\simeq 1,
\end{equation}
indicate a weak, large-scale, background-preserving deformation.

Thus, the viable SBC/F3 branch may be summarized as
\begin{equation}
\Lambda{\rm CDM\ background}
+
{\rm exponentially\ screened\ SBC/F3\ infrared\ perturbative\ sector}.
\end{equation}
